% EXP-004: OTel Weaver Registry to wasm4pm-compat Witness Generator
% Implementation Surface Chapter

\section{Implementation Surface}
\label{sec:otel-weaver-exp-004-registry-to-witness-implementation}

\subsection{Source Files}
\label{sec:otel-weaver-exp-004-registry-to-witness-source-files}

EXP-004 comprises three substantive source files and relies on two ecosystem-level files
for registry input and integration testing.

\textbf{\texttt{codegen/generator.rs}} is the executable Rust generator. It declares four
\texttt{serde}-derived structs for schema parsing (\texttt{ResolvedSchema},
\texttt{SchemaGroup}, \texttt{SchemaAttribute}), a \texttt{main} function that deserializes
the resolved schema JSON, filters for the \texttt{process.pi.activity} group, and writes
\texttt{generated\_witnesses.rs} to disk. Field name derivation uses Rust string replacement:
\texttt{attr.id.replace("process.pi.", "").replace(".", "\_")}. Attribute type mapping
converts \texttt{"string"} to \texttt{String} and \texttt{"int"} to \texttt{i64}, with
all other types defaulting to \texttt{String}. The generator also materializes the
\texttt{Lattice} trait definition and the full \texttt{PartialOrd} and \texttt{Lattice}
implementations for \texttt{ActivityWitness} as literal string pushes into the output
buffer.

The generator uses \texttt{std::fs::File::create} and \texttt{file.write\_all} to write
the output, creating parent directories via \texttt{fs::create\_dir\_all} if absent. The
\texttt{main} function signature is \texttt{fn main() -> Result<(), Box<dyn std::error::Error>>},
making all I/O failures explicit rather than panic-inducing.

\textbf{\texttt{src/generated\_witnesses.rs}} is the materialized output artifact. Its
header comments read:
\begin{verbatim}
//! Generated from OTel Weaver resolved registry. DO NOT EDIT HAND-CODED.
//! Generated at: 2026-06-01T10:10:51-07:00
\end{verbatim}
The file contains the \texttt{Lattice} trait, the \texttt{ActivityWitness} struct with
three public \texttt{String} fields carrying \texttt{serde rename} annotations, the
\texttt{PartialOrd} implementation, and the \texttt{Lattice} implementation. It exists on
disk and matches the README specification exactly.

\textbf{\texttt{README.md}} documents the full architectural pipeline as an ASCII art
diagram and contains the complete generator source listing as a Markdown code block.
Section~4 specifies the intended compilation command.

\textbf{\texttt{semconv/process\_pi.yaml}} (EXP-001) is the upstream source registry. It
declares \texttt{process.pi.activity} as a span group with eight attributes including
\texttt{process.pi.instance\_id}, \texttt{process.pi.activity.name},
\texttt{process.pi.lifecycle}, \texttt{process.pi.token.state\_before},
\texttt{process.pi.token.state\_after}, \texttt{process.pi.witness.id}, and
\texttt{process.pi.witness.hash}. The generator currently processes only the three
attributes inlined in its embedded JSON (instance\_id, witness.id, witness.hash); the
remaining five attributes from the full YAML are not reflected in the materialized output.

\textbf{\texttt{sources/wasm4pm/tests/weaver\_integration\_tests.rs}} provides the
integration test surface that exercises the complete feedstock-to-court pipeline including
\texttt{BridgeRx}, \texttt{validate\_and\_project}, \texttt{Admission}, and \texttt{Refusal}.

\subsection{Commands}
\label{sec:otel-weaver-exp-004-registry-to-witness-commands}

The README specifies the following compilation and execution sequence for the generator:
\begin{verbatim}
rustc --crate-type bin -O codegen/generator.rs -o codegen/generator \
    --edition 2021
./codegen/generator
\end{verbatim}
This command is architecturally incomplete as stated: \texttt{generator.rs} carries
\texttt{use serde::\{Deserialize, Serialize\}} and \texttt{use serde\_json} dependencies
that \texttt{rustc} cannot resolve without a \texttt{Cargo.toml} and compiled dependency
tree. The actual build path within the broader \texttt{wasm4pm} workspace is unverified.
This gap is documented in Section~\ref{sec:otel-weaver-exp-004-registry-to-witness-gaps}.

The integration tests are invoked via the standard Cargo test runner:
\begin{verbatim}
cargo test --test weaver_integration_tests
\end{verbatim}
The runtime checkpoint \texttt{GGEN\_OTEL\_WEAVER\_PI\_RUNTIME\_001.md} records the output
of this invocation: \texttt{test test\_weaver\_integration\_synthesis ... ok}, with all
62 tests across the ecosystem passing.

\subsection{Data and Ontology Surfaces}
\label{sec:otel-weaver-exp-004-registry-to-witness-data}

The convention registry \texttt{process\_pi.yaml} is the primary data surface. It is
authored in OTel Weaver YAML format (file format version 1.2.0, schema URL
\texttt{https://opentelemetry.io/schemas/1.25.0}) and defines the authoritative attribute
vocabulary for process-intelligence telemetry.

The mapping file \texttt{otel-weaver/mappings/otel-to-pi-witness-map.yaml} provides the
operational translation rules from OTel resource attributes to process-intelligence
feedstock keys. It defines three witness mapping rules: service identity to
\texttt{pi.feedstock.witness.id} (via format string \texttt{wit\_\{service.namespace\}\_\{service.name\}}),
SDK identity to \texttt{pi.feedstock.source}, and payload size coalescing. The file also
declares three witness registry entries for eBPF system call, database audit, and API
gateway contexts. Its conformance isolation rule states: ``Never map feedstock directly to
court consequences. A witness only asserts observations of feedstock, not conformity
verdicts.''

The generator's inline JSON schema is a subset of the full \texttt{process\_pi.yaml}
definition, containing only the three attributes (\texttt{process.pi.instance\_id},
\texttt{process.pi.witness.id}, \texttt{process.pi.witness.hash}) that the generator
currently materializes. This inlining is the primary evidence for the open question about
whether genuine Weaver CLI resolution was ever executed against the full registry.

\subsection{Templates and Rendered Artifacts}
\label{sec:otel-weaver-exp-004-registry-to-witness-templates}

The generator does not use an external template engine. Template logic is embedded directly
in the \texttt{generator.rs} source as \texttt{String::push\_str} calls and
\texttt{format!} macro invocations. The output structure is fixed: header comment block,
\texttt{use} declarations, \texttt{Lattice} trait definition, one \texttt{ActivityWitness}
struct per matching group, and one \texttt{Lattice} implementation block per struct.

The rendered artifact \texttt{src/generated\_witnesses.rs} carries the instruction
\texttt{DO NOT EDIT HAND-CODED} in its header, establishing the generator as the sole
authoritative source for its content. Any hand-edit to the materialized file violates the
provenance contract implied by this header and should be treated as a manufacturing defect.

The \texttt{top()} element for \texttt{ActivityWitness} is materialized as:
\begin{verbatim}
witness_hash: "f".repeat(64)
\end{verbatim}
This produces a 64-character string of the character \texttt{f}, representing the maximal
BLAKE3 hex value. This encoding is a convention chosen by the generator author and is not
derived from any external specification; it is noted as a theoretical approximation
requiring formal justification (see Section~\ref{sec:otel-weaver-exp-004-registry-to-witness-gaps}).
